% ===========================================================================
%  Bagian 3 -- Benchmark
% ===========================================================================
\section{Benchmark}\label{sec:benchmark}

Seluruh angka pada bagian ini dihitung ulang oleh
\code{scripts/01\_tables.R} dan \code{scripts/02\_figures.R} langsung dari
sembilan berkas \code{inst/extdata/*.rds}. Tidak ada angka yang diketik tangan,
dan setiap tabel/figur mencantumkan file sumbernya pada caption. Bila klaim di
\code{README.Rmd} atau vignette bertolak belakang dengan file data, laporan ini
memakai data dan mencatat kontradiksinya pada bagian~\ref{sec:kontradiksi}
(daftar lengkap: \code{DISCREPANCIES.md} bagian H).

\subsection{Metodologi}\label{sec:metodologi-bench}

\paragraph{Apa yang diukur.}
Semua file benchmark berformat sama: kolom \code{Method},
\code{Min (s)}, \code{Median (s)}, \code{Iterations/sec},
\code{Memory (MB)}, \code{n\_iter}, \code{n}. Dua kolom waktu berasal dari
\code{bench::mark()\$min} dan \code{\$median}; \code{Memory (MB)} =
\code{as.numeric(res\$mem\_alloc)/1024\textasciicircum2}, yaitu \emph{total alokasi memori
selama run}, bukan puncak RSS (lihat \code{benchmarks/bench\_beta\_tipsae.R:92-98}).
\code{n\_iter} mencatat jumlah iterasi aktual yang benar-benar dijalankan.

\paragraph{Skrip, pembanding, ukuran, iterasi, seed.}
Tabel~\ref{tab:skrip-bench} disusun dari pembacaan langsung 11 skrip di
\code{benchmarks/}.

\begin{table}[htbp]
\centering
\caption{Skrip benchmark: pembanding, ukuran masalah, iterasi, dan seed. Sumber:
pembacaan \code{benchmarks/*.R} (lihat \code{\_notes/bench\_methodology.md}).}
\label{tab:skrip-bench}
\small
\footnotesize
\begin{tabularx}{\textwidth}{@{}>{\raggedright\arraybackslash}p{4.9cm}
  >{\raggedright\arraybackslash}X
  l
  >{\raggedright\arraybackslash}p{4.6cm}@{}}
\toprule
Skrip & Pembanding (dari kode) & $n$ & Iterasi / seed \\
\midrule
\code{bench\_beta\_tipsae.R} &
  \code{fastsae::hb\_area} vs \code{tipsae::fit\_sae} &
  250, 500 & 20; data \code{seed = 42+n}; tipsae \code{seed = 42} \\
\code{bench\_spatial\_beta.R} &
  \code{fastsae} vs \code{tipsae} (+ \code{sf}, \code{spdep}) &
  30, 50, 100 & 20; data \code{42+n} \\
\code{bench\_spatio\_temporal\_beta.R} &
  \code{fastsae} vs \code{tipsae} ($T=5$) &
  30\,--\,1000 & 5; data \code{42+n} \\
\code{bench\_hb\_normal.R} &
  \code{fastsae::hb\_area} vs \code{hbsae::fSAE.Area} &
  30, 50, 100, 250 & 10; data \code{42+n} \\
\code{update\_hb\_benchmark.R} &
  \code{fastsae} vs \code{hbsae} &
  500, 1000 & 10; data \code{42+n} \\
\code{run\_hb\_area\_full\_benchmark.R} &
  \code{fastsae} vs \code{fastsaehb} vs \code{tipsae} (Stan HMC) &
  skenario tunggal & \code{set.seed(123/42)}; tidak menulis file \\
\code{run\_\allowbreak{}multipackage\_\allowbreak{}hb\_\allowbreak{}benchmark.R} &
  \code{fastsae} vs \code{hbsae}, \code{CARBayes}, \code{fastsaehb},
  \code{tipsae} &
  --- & 5 iterasi; tidak menulis file \\
\code{compare\_ebp\_tipsae\_truth.R},
\code{sim\_finite\_population\_beta.R},
\code{sim\_mc\_finite\_population.R},
\code{sim\_mc\_finite\_pop\_d500.R} &
  \code{fastsae} vs \code{tipsae} vs estimasi langsung (akurasi terhadap
  kebenaran model) &
  $D=60$/$500$ &
  $R=100$ (seed 2024) dan $R=25$ (seed 999); \code{mclapply} 5--6 core \\
\bottomrule
\end{tabularx}
\end{table}

\paragraph{Apa yang tidak tercatat.}
Model CPU, jumlah inti, RAM, sistem operasi, tanggal run, versi \pkg{R} pada
saat run, dan versi paket pembanding \emph{tidak ada} di dalam file
\code{.rds}. Satu-satunya pernyataan hardware di repo adalah
\code{paper/article.tex:277} (``Apple Silicon processor under R 4.4.2'').
Tidak ada skrip yang mengatur \code{OMP\_NUM\_THREADS}; \code{README.Rmd:126}
hanya menulis bahwa ``timings depend on hardware and the number of threads''.
Empat file (\code{eblup\_benchmark.rds}, \code{seblup\_benchmark.rds},
\code{steblup\_benchmark.rds}, \code{beta\_estimates\_n1000.rds}) \emph{tidak
punya skrip pembuat} di repo mana pun (terverifikasi dengan
\code{git log -S}); karenanya hasilnya tidak dapat direproduksi dari repo.

\paragraph{Keadilan perbandingan.}
Pada skrip yang ada, kedua metode menerima \emph{data simulasi yang sama}
(\code{sim\_area\_data(D = n\_val, seed = 42 + n\_val)}), formula yang sama,
dan kolom \code{vardir} yang sama; \code{bench::mark(..., check = FALSE,
filter\_gc = FALSE, memory = TRUE, time\_unit = "s")}. Perbedaan yang tetap
ada dan harus dibaca sebagai \emph{caveat}:

\begin{itemize}
  \item \code{tipsae} dijalankan sebagai pemanggilan MCMC/Stan dengan
        \code{seed = 42} dan \code{refresh = 0}, sedangkan \pkg{fastsae}
        memakai aproksimasi Laplace INLA --- dua algoritme yang berbeda
        secara mendasar, bukan dua implementasi algoritme yang sama.
  \item \code{hbsae::fSAE.Area} adalah pemanggilan langsung tanpa
        penyetelan iterasi; \code{fastsae::hb\_area} memakai default
        \code{strategy = "laplace"}.
  \item Untuk tiga file pertama (\code{eblup}, \code{seblup},
        \code{steblup}) skrip pembuatnya tidak ada, sehingga kesetaraan
        pengaturan awal, kriteria konvergensi, dan jumlah thread tidak dapat
        diverifikasi dari repo.
  \item Beberapa \code{n\_iter} aktual berbeda dari nilai yang diminta
        (lihat kolom \code{n\_iter} di file), tanda run besar memakai argumen
        \code{--iter} yang tidak tercatat di file mana pun.
\end{itemize}

\subsection{Waktu komputasi}\label{sec:waktu}

Tabel~\ref{tab:time-fh}--\ref{tab:time-stfh} memuat waktu median per iterasi
untuk tiga keluarga Fay--Herriot frekuentis; Tabel~\ref{tab:time-beta}
sampai Tabel~\ref{tab:time-hb} untuk model beta dan model hierarki Bayes.
Pembandingnya: \pkg{sae} dan \pkg{emdi} untuk FH, \pkg{tipsae} untuk beta,
\pkg{hbsae} untuk Bayes Gaussian.

Perlu ditegaskan apa yang dimaksud ``keluarga beta'' di sini: \pkg{fastsae}
tidak mengekspor fungsi khusus beta --- yang dijalankan adalah
\code{fastsae::hb\_area} dengan \code{family = "beta"}, sedangkan
pembandingnya \code{tipsae::fit\_sae} (\code{benchmarks/bench\_beta\_tipsae.R:9-10}).
Rumus regresi beta yang mendasari pembanding itu dibahas
\citet{ferrari2004beta} dan \citet{janicki2020beta}, perluasannya ke SAE
area-level oleh \citet{esteban2019compositional} dan ke campuran beta oleh
\citet{denicolo2024betamixture}; rujukan paketnya \citet{denicolo2024tipsae}.

\begin{table}[htbp]
\centering
\caption{Waktu median per iterasi (detik) menurut jumlah area $n$. Sumber: file \texttt{eblup\_benchmark.rds} di \texttt{inst/extdata} (repo \texttt{fastsae}); angka dihitung ulang oleh \texttt{scripts/01\_tables.R}, ditampilkan hingga 6 digit signifikan.}
\label{tab:time-fh}
\small
\begin{tabular}{lrrrr}
\toprule
$n$ & fastsae & emdi & sae \\
\midrule
30 & 0.000588063 & 0.0108802 & 0.00155167 \\
50 & 0.000510983 & 0.0180643 & 0.0019507 \\
100 & 0.00055022 & 0.0505076 & 0.00344267 \\
250 & 0.000899663 & 0.857347 & 0.037575 \\
500 & 0.00140398 & 5.86366 & 0.196726 \\
1000 & 0.00413358 & 51.0224 & 1.50363 \\
\bottomrule
\end{tabular}
\end{table}


\begin{table}[htbp]
\centering
\caption{Waktu median per iterasi (detik) menurut jumlah area $n$. Sumber: file \texttt{seblup\_benchmark.rds} di \texttt{inst/extdata} (repo \texttt{fastsae}); angka dihitung ulang oleh \texttt{scripts/01\_tables.R}, ditampilkan hingga 6 digit signifikan.}
\label{tab:time-sfh}
\small
\begin{tabular}{lrrrr}
\toprule
$n$ & fastsae & emdi & sae \\
\midrule
30 & 0.000734618 & 0.0107105 & 0.00410111 \\
50 & 0.00114812 & 0.0180005 & 0.00967977 \\
100 & 0.00472576 & 0.05299 & 0.0705356 \\
250 & 0.0238317 & 0.843782 & 1.07476 \\
500 & 0.142051 & 3.13662 & 7.11793 \\
1000 & 0.832158 & 51.0593 & 67.1518 \\
\bottomrule
\end{tabular}
\end{table}


\begin{table}[htbp]
\centering
\caption{Waktu median per iterasi (detik) menurut jumlah area $n$. Model spasio-temporal dengan $T$ periode waktu; waktu termasuk estimasi komponen varian. Sumber: file \texttt{steblup\_benchmark.rds} di \texttt{inst/extdata} (repo \texttt{fastsae}); angka dihitung ulang oleh \texttt{scripts/01\_tables.R}, ditampilkan hingga 6 digit signifikan.}
\label{tab:time-stfh}
\small
\begin{tabular}{lrrr}
\toprule
$n$ & fastsae & sae \\
\midrule
30 & 0.0367356 & 0.118558 \\
50 & 0.0898305 & 0.420018 \\
100 & 2.29928 & 17.9425 \\
250 & 3.81651 & 58.4315 \\
500 & 10.3596 & 253.343 \\
1000 & 58.1005 & 1786.49 \\
\bottomrule
\end{tabular}
\end{table}



\begin{table}[htbp]
\centering
\caption{Waktu median per iterasi (detik) menurut jumlah area $n$. Model area-level beta (\texttt{fastsae::hb\_area} vs \texttt{tipsae::fit\_sae}). Sumber: file \texttt{beta\_benchmark.rds} di \texttt{inst/extdata} (repo \texttt{fastsae}); angka dihitung ulang oleh \texttt{scripts/01\_tables.R}, ditampilkan hingga 6 digit signifikan.}
\label{tab:time-beta}
\small
\begin{tabular}{lrrr}
\toprule
$n$ & fastsae & tipsae \\
\midrule
30 & 1.24489 & 2.34867 \\
50 & 1.33152 & 3.77451 \\
100 & 1.27232 & 8.23587 \\
250 & 1.31454 & 26.2762 \\
500 & 1.41239 & 64.2141 \\
1000 & 1.63967 & 187.162 \\
\bottomrule
\end{tabular}
\end{table}


\begin{table}[htbp]
\centering
\caption{Waktu median per iterasi (detik) menurut jumlah area $n$. Model area-level beta dengan efek spasial. Sumber: file \texttt{spatial\_beta\_benchmark.rds} di \texttt{inst/extdata} (repo \texttt{fastsae}); angka dihitung ulang oleh \texttt{scripts/01\_tables.R}, ditampilkan hingga 6 digit signifikan.}
\label{tab:time-sbeta}
\small
\begin{tabular}{lrrr}
\toprule
$n$ & fastsae & tipsae \\
\midrule
30 & 0.613834 & 25.7233 \\
50 & 0.620765 & 65.7235 \\
100 & 1.37316 &  100.02 \\
250 & 1.31548 & 155.898 \\
500 & 1.43413 & 248.821 \\
1000 & 1.72812 & 404.342 \\
\bottomrule
\end{tabular}
\end{table}


\begin{table}[htbp]
\centering
\caption{Waktu median per iterasi (detik) menurut jumlah area $n$. Model beta spasio-temporal ($T=5$); data tersedia hanya untuk $n=30$ dan $n=50$. Sumber: file \texttt{spatio\_temporal\_beta\_benchmark.rds} di \texttt{inst/extdata} (repo \texttt{fastsae}); angka dihitung ulang oleh \texttt{scripts/01\_tables.R}, ditampilkan hingga 6 digit signifikan.}
\label{tab:time-stbeta}
\small
\begin{tabular}{lrrr}
\toprule
$n$ & fastsae & tipsae \\
\midrule
30 & 1.44206 & 540.573 \\
50 & 1.66767 & 908.173 \\
\bottomrule
\end{tabular}
\end{table}


\begin{table}[htbp]
\centering
\caption{Waktu median per iterasi (detik) menurut jumlah area $n$. Model hierarki Bayes area-level Gaussian (\texttt{hb\_area}) vs \texttt{hbsae::fSAE.Area}. Sumber: file \texttt{hb\_benchmark.rds} di \texttt{inst/extdata} (repo \texttt{fastsae}); angka dihitung ulang oleh \texttt{scripts/01\_tables.R}, ditampilkan hingga 6 digit signifikan.}
\label{tab:time-hb}
\small
\begin{tabular}{lrrr}
\toprule
$n$ & fastsae & hbsae \\
\midrule
30 & 1.29238 & 0.0209383 \\
50 & 1.28063 & 0.0255576 \\
100 &  1.2507 & 0.0373128 \\
250 & 1.22658 & 0.0744251 \\
500 & 1.25466 & 0.143842 \\
1000 & 1.27669 & 0.271585 \\
\bottomrule
\end{tabular}
\end{table}



Tiga pengamatan yang terbaca langsung dari tabel:

\begin{enumerate}[label=\arabic*.]
  \item Pada masalah kecil ($n=30$) selisih antar paket kecil bahkan untuk
        FH, sebab biaya tetap (parsing, pembangunan matriks desain) masih
        dominan.
  \item Pada FH area-level, waktu pembanding tumbuh jauh lebih cepat
        daripada \pkg{fastsae} seiring $n$; pada $n$ terbesar pada file,
        \pkg{emdi} membutuhkan waktu dua orde magnitudo lebih besar.
  \item Pada model Bayes Gaussian justru \pkg{hbsae} yang lebih cepat
        di seluruh $n$ pada file --- tidak ada pembanding lain pada skenario
        ini, dan temuan ini dibahas pada bagian~\ref{sec:kaveat}.
\end{enumerate}

\subsection{Penggunaan memori}\label{sec:memori-bench}

\begin{table}[htbp]
\centering
\caption{Alokasi memori total per iterasi (MB) menurut jumlah area $n$. Sumber: file \texttt{eblup\_benchmark.rds} di \texttt{inst/extdata} (repo \texttt{fastsae}); angka dihitung ulang oleh \texttt{scripts/01\_tables.R}, ditampilkan hingga 6 digit signifikan.}
\label{tab:mem-fh}
\small
\begin{tabular}{lrrrr}
\toprule
$n$ & fastsae & emdi & sae \\
\midrule
30 & 0.00506592 & 4.18307 & 0.153381 \\
50 & 0.00884247 & 10.4676 & 0.373138 \\
100 & 0.0162811 & 37.3394 & 0.917542 \\
250 & 0.0385971 & 240.893 & 6.34988 \\
500 & 0.0757904 & 880.314 & 18.2391 \\
1000 & 0.185699 &  3773.6 & 71.9765 \\
\bottomrule
\end{tabular}
\end{table}


\begin{table}[htbp]
\centering
\caption{Alokasi memori total per iterasi (MB) menurut jumlah area $n$. Sumber: file \texttt{seblup\_benchmark.rds} di \texttt{inst/extdata} (repo \texttt{fastsae}); angka dihitung ulang oleh \texttt{scripts/01\_tables.R}, ditampilkan hingga 6 digit signifikan.}
\label{tab:mem-sfh}
\small
\begin{tabular}{lrrrr}
\toprule
$n$ & fastsae & emdi & sae \\
\midrule
30 & 0.029808 & 4.18307 & 1.37314 \\
50 & 0.0737839 & 10.4676 & 3.31219 \\
100 & 0.258606 & 37.3394 & 17.9755 \\
250 & 1.49972 & 240.893 & 110.388 \\
500 & 5.85706 & 949.864 & 776.407 \\
1000 & 23.1548 & 3772.57 & 1920.01 \\
\bottomrule
\end{tabular}
\end{table}


\begin{table}[htbp]
\centering
\caption{Alokasi memori total per iterasi (MB) menurut jumlah area $n$. Sumber: file \texttt{steblup\_benchmark.rds} di \texttt{inst/extdata} (repo \texttt{fastsae}); angka dihitung ulang oleh \texttt{scripts/01\_tables.R}, ditampilkan hingga 6 digit signifikan.}
\label{tab:mem-stfh}
\small
\begin{tabular}{lrrr}
\toprule
$n$ & fastsae & sae \\
\midrule
30 & 0.0274048 & 43.8007 \\
50 & 0.0466537 & 107.829 \\
100 & 0.0901871 &  2551.1 \\
250 & 0.225403 & 3624.71 \\
500 & 0.448158 & 8124.37 \\
1000 & 0.893669 & 32478.3 \\
\bottomrule
\end{tabular}
\end{table}


\begin{table}[htbp]
\centering
\caption{Alokasi memori total per iterasi (MB) menurut jumlah area $n$. Model hierarki Bayes area-level Gaussian. Sumber: file \texttt{hb\_benchmark.rds} di \texttt{inst/extdata} (repo \texttt{fastsae}); angka dihitung ulang oleh \texttt{scripts/01\_tables.R}, ditampilkan hingga 6 digit signifikan.}
\label{tab:mem-hb}
\small
\begin{tabular}{lrrr}
\toprule
$n$ & fastsae & hbsae \\
\midrule
30 & 46.7163 & 5.93539 \\
50 & 2.40248 &  5.8022 \\
100 & 3.59152 & 12.3786 \\
250 & 7.15878 & 45.3857 \\
500 & 57.7934 & 112.206 \\
1000 & 25.0034 & 220.352 \\
\bottomrule
\end{tabular}
\end{table}


\begin{table}[htbp]
\centering
\caption{Ringkasan alokasi memori total per iterasi (MB): rerata, minimum dan maksimum lintas seluruh $n$ pada tiap file sumber. Kolom ``maks'' adalah nilai pada $n$ terbesar pada file tersebut.}
\label{tab:mem-summary}
\small
\begin{tabular}{llrrr}
\toprule
File & Metode & Rerata & Min & Maks \\
\midrule
\texttt{eblup} & fastsae & 0.0550461 & 0.00506592 & 0.185699 \\
\texttt{eblup} & emdi & 824.466 & 4.18307 &  3773.6 \\
\texttt{eblup} & sae & 16.3349 & 0.153381 & 71.9765 \\
\texttt{seblup} & fastsae & 5.14563 & 0.029808 & 23.1548 \\
\texttt{seblup} & emdi & 835.886 & 4.18307 & 3772.57 \\
\texttt{seblup} & sae & 471.577 & 1.37314 & 1920.01 \\
\texttt{steblup} & fastsae & 0.288579 & 0.0274048 & 0.893669 \\
\texttt{steblup} & sae & 7821.69 & 43.8007 & 32478.3 \\
\texttt{beta} & fastsae & 81.0359 & 2.59546 & 133.585 \\
\texttt{beta} & tipsae & 5860.87 &  499.97 & 18335.2 \\
\texttt{spatial\_beta} & fastsae & 70.6704 & 3.36058 & 161.025 \\
\texttt{spatial\_beta} & tipsae & 7366.18 & 642.835 & 23038.7 \\
\texttt{spatio\_temporal\_beta} & fastsae & 105.869 &  102.73 & 109.008 \\
\texttt{spatio\_temporal\_beta} & tipsae &  3913.6 & 2947.08 & 4880.13 \\
\texttt{hb} & fastsae & 23.7777 & 2.40248 & 57.7934 \\
\texttt{hb} & hbsae &   67.01 &  5.8022 & 220.352 \\
\bottomrule
\end{tabular}
\end{table}



Kolom \code{Memory (MB)} adalah total alokasi selama run (bagian
\ref{sec:metodologi-bench}), sehingga angkanya lebih besar daripada puncak RSS
yang dilaporkan \code{bench::mark()} sebagai \code{mem\_usage}. Klaim
\code{README.Rmd:124} (``$\mathtt{<10}$ MB'' untuk \pkg{fastsae}) dan label
``Peak Memory'' terbantahkan oleh file data pada beberapa baris; rinciannya di
bagian~\ref{sec:kontradiksi}.

\subsection{Skala terhadap jumlah area}\label{sec:skala}

Gambar~\ref{fig:time-fh} dan Gambar~\ref{fig:time-beta} menampilkan waktu
median terhadap $n$ pada skala log--log untuk seluruh keluarga model;
Gambar~\ref{fig:mem} menampilkan memori; Gambar~\ref{fig:ratio} merangkum
rasio waktu pembanding/tes terhadap $n$.

\begin{figure}[htbp]
  \centering
  \includegraphics[width=\textwidth]{fig-time-fh.pdf}
  \caption{Waktu median per iterasi terhadap jumlah area $n$ (skala log pada
  kedua sumbu). Sumber: \texttt{inst/extdata/eblup\_benchmark.rds},
  \texttt{seblup\_benchmark.rds}, \texttt{steblup\_benchmark.rds},
  \texttt{hb\_benchmark.rds}; digambar oleh \texttt{scripts/02\_figures.R}.}
  \label{fig:time-fh}
\end{figure}

\begin{figure}[htbp]
  \centering
  \includegraphics[width=\textwidth]{fig-time-beta.pdf}
  \caption{Waktu median untuk tiga varian model beta (skala log--log).
  Sumber: \texttt{inst/extdata/beta\_benchmark.rds},
  \texttt{spatial\_beta\_benchmark.rds},
  \texttt{spatio\_temporal\_beta\_benchmark.rds}; \texttt{scripts/02\_figures.R}.
  Panel (c) hanya memiliki $n=30$ dan $n=50$.}
  \label{fig:time-beta}
\end{figure}

\begin{figure}[htbp]
  \centering
  \includegraphics[width=\textwidth]{fig-mem.pdf}
  \caption{Alokasi memori total per iterasi terhadap $n$ (skala log--log).
  Sumber: file \texttt{inst/extdata/*\_benchmark.rds} sebagaimana tercantum pada
  tiap panel; \texttt{scripts/02\_figures.R}.}
  \label{fig:mem}
\end{figure}

\begin{figure}[htbp]
  \centering
  \includegraphics[width=\textwidth]{fig-ratio.pdf}
  \caption{Rasio waktu (pembanding/tes) terhadap $n$; garis putus-putus
  menandai rasio $=1$. Sumber: file \texttt{inst/extdata/*\_benchmark.rds};
  \texttt{scripts/02\_figures.R}.}
  \label{fig:ratio}
\end{figure}

Tabel~\ref{tab:ratio} merangkum seluruh rasio per $n$. Rentangnya
jauh lebih lebar daripada klaim di \code{README.Rmd}: pada ujung bawah ada
skenario tempat pembanding lebih cepat dari \pkg{fastsae}, dan pada ujung
atas rasio mencapai empat orde magnitudo --- karenanya laporan memilih
menampilkan rentang, bukan satu angka tunggal.

\begin{table}[htbp]
\centering
\caption{Rasio kecepatan $=$ median waktu pembanding $/$ median waktu \texttt{fastsae} (semakin besar semakin cepat \texttt{fastsae}; nilai $<1$ berarti pembanding lebih cepat). Rasio dihitung per $n$ lalu diringkas.}
\label{tab:ratio}
\small
\begin{tabular}{llrrrr}
\toprule
Model & Pembanding & $n$ pertama & $n$ terbesar & min & max \\
\midrule
FH area-level (\texttt{eblup\_fh}) & emdi & 18.5018 & 12343.4 & 18.5018 & 12343.4 \\
 & \multicolumn{5}{r}{\emph{rata-rata lintas $n$: 2936.41$\times$}} \\[-2pt]
FH area-level (\texttt{eblup\_fh}) & sae &  2.6386 &  363.76 &  2.6386 &  363.76 \\
 & \multicolumn{5}{r}{\emph{rata-rata lintas $n$: 93.0598$\times$}} \\[-2pt]
FH spasial (\texttt{eblup\_sfh}) & emdi & 14.5797 & 61.3577 &  11.213 & 61.3577 \\
 & \multicolumn{5}{r}{\emph{rata-rata lintas $n$: 26.7192$\times$}} \\[-2pt]
FH spasial (\texttt{eblup\_sfh}) & sae & 5.58264 &  80.696 & 5.58264 &  80.696 \\
 & \multicolumn{5}{r}{\emph{rata-rata lintas $n$: 34.1403$\times$}} \\[-2pt]
FH spasio-temporal (\texttt{eblup\_stfh}) & sae & 3.22733 & 30.7484 & 3.22733 & 30.7484 \\
 & \multicolumn{5}{r}{\emph{rata-rata lintas $n$:   14.37$\times$}} \\[-2pt]
Beta area-level (\texttt{hb\_area}) & tipsae & 1.88664 & 114.146 & 1.88664 & 114.146 \\
 & \multicolumn{5}{r}{\emph{rata-rata lintas $n$: 31.7991$\times$}} \\[-2pt]
Beta spasial & tipsae &  41.906 & 233.978 &  41.906 & 233.978 \\
 & \multicolumn{5}{r}{\emph{rata-rata lintas $n$: 124.435$\times$}} \\[-2pt]
Beta spasio-temporal & tipsae & 374.863 & 544.576 & 374.863 & 544.576 \\
 & \multicolumn{5}{r}{\emph{rata-rata lintas $n$:  459.72$\times$}} \\[-2pt]
Hierarki Bayes Gaussian & hbsae & 0.0162013 & 0.212725 & 0.0162013 & 0.212725 \\
 & \multicolumn{5}{r}{\emph{rata-rata lintas $n$: 0.0756733$\times$}} \\[-2pt]
\bottomrule
\end{tabular}
\end{table}



\subsection{Ekuivalensi numerik terhadap pembanding}\label{sec:ekuivalensi}

File \code{beta\_estimates\_n1000.rds} memuat pasangan estimasi dan MSE untuk
1000 area antara \pkg{fastsae} dan \pkg{tipsae} pada model beta yang sama.
Tabel~\ref{tab:equiv} menghitung metrik kesesuaian dari pasangan kolom itu;
Gambar~\ref{fig:equiv} menampilkannya secara grafis.

\begin{table}[htbp]
\centering
\caption{Ekuivalensi numerik antara \texttt{fastsae} dan \texttt{tipsae} pada $n=1000$ area (indikator beta). Semua metrik dihitung dari pasangan kolom pada file sumber.}
\label{tab:equiv}
\small
\begin{tabular}{lr}
\toprule
Metrik & Nilai \\
\midrule
Pearson $r$ (estimasi) & 0.999343 \\
Spearman $\rho$ (estimasi) & 0.999208 \\
MAE (estimasi) & 0.0196916 \\
RMSD (estimasi) & 0.0236901 \\
Galat absolut maksimum & 0.0558335 \\
Pearson $r$ (MSE) & 0.963148 \\
Median selisih mutlak MSE & 0.000171686 \\
Median rasio MSE fastsae/tipsae & 2.21687 \\
\bottomrule
\end{tabular}
\end{table}



\begin{figure}[htbp]
  \centering
  \includegraphics[width=\textwidth]{fig-equiv.pdf}
  \caption{Ekuivalensi numerik \texttt{fastsae} vs \texttt{tipsae} untuk
  $n=1000$ area: (a) prediktor area terhadap garis identitas;
  (b) MSE pada skala log. Sumber: \texttt{inst/extdata/beta\_estimates\_n1000.rds};
  \texttt{scripts/02\_figures.R}.}
  \label{fig:equiv}
\end{figure}

Kedua paket memprediksi indikator yang sama dengan selisih kecil pada prediktor
(korelasi sangat tinggi; Tabel~\ref{tab:equiv}), sedangkan MSE berbeda lebih
nyata karena kedua paket memakai pendekatan ketidakpastian yang berbeda (INLA
Laplace vs posterior MCMC). Angka MSE tidak boleh dibaca sebagai ``satu paket
lebih benar'' tanpa rujukan kebenaran; pembanding terhadap kebenaran model
tersedia pada skrip \code{benchmarks/compare\_ebp\_tipsae\_truth.R} dan
\code{sim\_mc\_finite\_population.R} (hasil tersimpan di
\code{benchmarks/output/}, di luar cakupan \code{inst/extdata}).

\subsection{Kaveat dan keadilan}\label{sec:kaveat}

\begin{enumerate}[label=\arabic*.]
  \item \textbf{Beda algoritme, bukan beda implementasi.} \pkg{tipsae} adalah
        Bayes MCMC/Stan \citep{denicolo2024tipsae}, \pkg{hbsae} memakai
        pendekatan Bayesnya sendiri, \pkg{emdi} dan \pkg{sae} memakai
        \pkg{lme4}/aljabar \pkg{Matrix} di R, sedangkan \pkg{fastsae} memakai
        C++ + INLA. Selisih waktu mencerminkan perbedaan ini sekaligus.
  \item \textbf{Tanpa penetapan thread.} Tidak ada skrip yang mengunci
        \code{OMP\_NUM\_THREADS}; bootstrap paralel (bagian
        \ref{sec:bootstrap-paralel}) dapat memakai seluruh core mesin,
        sedangkan pembanding berjalan serial. Ini membuat perbandingan
        \emph{memihak} ke \pkg{fastsae} pada skala besar, dan dilaporkan apa
        adanya.
  \item \textbf{Kriteria konvergensi tidak disamakan.} \pkg{fastsae} memakai
        toleransi $10^{-4}$ dan 100 iterasi (Tabel~\ref{tab:konvergensi});
        pengaturan pembanding tidak tercatat di file.
  \item \textbf{Memori bukan puncak.} Kolom memori adalah akumulasi alokasi,
        sehingga tidak setara dengan \emph{peak RSS} yang lazim dilaporkan.
  \item \textbf{Empat file tanpa skrip.} Hasil FH (tiga file) dan
        \code{beta\_estimates\_n1000} tidak dapat direproduksi dari repo;
        laporan memakainya apa adanya dan menandainya.
  \item \textbf{Run inkremental.} Skrip 1--5 menulis file secara inkremental
        (baris lama per $n$ diganti), sehingga file bisa berisi campuran run
        dari waktu yang berbeda.
\end{enumerate}

\subsection{Kontradiksi klaim dokumentasi terhadap data}\label{sec:kontradiksi}

Aturan laporan: data menang. Temuan berikut berasal dari perbandingan langsung
klaim \code{README.Rmd}/\code{vignettes/benchmarks.Rmd}/\code{paper/article.tex}
dengan isi \code{inst/extdata} (perhitungan lengkap:
\code{\_notes/bench\_methodology.md} bagian 6 dan 8).

\begin{table}[htbp]
\centering
\caption{Klaim dokumentasi yang bertolak belakang dengan file data. Sumber:
\texttt{README.Rmd}, \texttt{vignettes/benchmarks.Rmd}, \texttt{paper/article.tex}
vs \texttt{inst/extdata/*.rds} (butir DOC-1\,--\,DOC-11 di \texttt{DISCREPANCIES.md}).}
\label{tab:kontradiksi}
\small
\begin{tabularx}{\textwidth}{@{}l X X@{}}
\toprule
Butir & Klaim dokumentasi & Fakta dari data \\
\midrule
DOC-1 & \code{README.Rmd:118} menandai tabel ``Benchmark across $n=1000$
  areas'' & angka tabel adalah rata-rata lintas $n=30\ldots1000$ \\
DOC-2 & \code{README.Rmd:122} ``$\sim$190$\times$'' dan ``$\sim$6400$\times$'' &
  rasio yang benar tidak cocok dengan salah satunya (lihat Tabel~\ref{tab:ratio}) \\
DOC-3 & \code{README.Rmd:124} fastsae ``$<$ 10 MB'' &
  beberapa baris \pkg{fastsae} melebihi 10 MB (Tabel~\ref{tab:mem-summary}) \\
DOC-4 & \code{README.Rmd:124} pembanding ``$\sim$400 MB''/``$\sim$800 MB'' &
  campuran dua file; untuk FH angkanya jauh di bawah itu \\
DOC-5/6 & \code{README.Rmd:34-35} ``50 to 6{,}000$\times$'' dan
  \code{:87} ``10--100$\times$'' &
  rentang teramati lebih lebar di kedua ujung (Tabel~\ref{tab:ratio}) \\
DOC-7 & \code{vignettes/benchmarks.Rmd:39} emdi spasial ``8{,}69 s'' &
  nilai pada file berbeda (Tabel~\ref{tab:time-sfh}) \\
DOC-8 & label ``\emph{Peak} Memory'' di README, vignette, dan artikel &
  nilai adalah rerata \code{mem\_alloc}, bukan puncak \\
DOC-9 & \code{paper/article.tex:100,480} ``73$\times$ to 544$\times$'' &
  ada nilai serendah $1{,}887\times$ (Tabel~\ref{tab:ratio}) \\
DOC-10 & \code{vignettes/}\allowbreak{}\code{benchmarks.Rmd:407-414,}\allowbreak{}\code{634-639} kolom Spatial/ST
  $n\ge100$ & tidak ada data pendukung di \code{inst/extdata} \\
DOC-11 & reproduksi hasil & empat file tidak punya skrip pembuat di repo \\
\bottomrule
\end{tabularx}
\end{table}

Implikasinya langsung: laporan ini tidak mengutip satu pun angka kecepatan
dari README. Angka yang dipakai seluruhnya berasal dari tabel dan gambar di
atas, yang kesemuanya dihasilkan ulang dari file data.
